\section{The full primary theta sum and its integral lattice index}
\label{sec:geometry-primary-trace}

The formal trace of the preceding section extends to the full primary
torsion group with the original stabilizer factor.  A separate local
lattice construction identifies that factor with the primary part of
the denominator.  The regularity of the original translated sections
is retained explicitly; it is not a consequence of rational torsion.

\subsection*{The actual adapted lattice}

Let $E/F$ be the CM extension in the original host, unramified at every
finite prime, and let $\Lambda\subset E^n$ be self-dual for the standard
Hermitian form.  Write $\Lambda_0=\mathcal O_E^n$ and
$\Lambda_\cap=\Lambda\cap\Lambda_0$.  Suppose $N$ is positive and odd and
$N\Lambda_0\subset\Lambda_\cap$.  For a rational prime $\ell$, put
$S=\{\mathfrak p:\mathfrak p\mid2\ell\}$ in $\mathcal O_E$.

\begin{proposition}\label{prop:primary-adapted-construction}
There are positive $k,k_1$ prime to $2\ell$, a unitary
$\gamma\in U_n(E/F)$ and matrices $p,q,g\in M_n(\mathcal O_E)$ such that
\begin{align}
 pq&=kN\gamma,&gq&=k_1N I_n,\\
 k\gamma\Lambda&\subset\Lambda_0,&
 k\Lambda_0&\subset\gamma\Lambda,\\
 g\Lambda_0&\subset\Lambda_\cap,&
 k_1\Lambda_\cap&\subset g\Lambda_0.
 \label{eq:primary-adapted-data}
\end{align}
In particular $(pq)^*(pq)=(kN)^2I_n$.  For every $r\geq0$, setting
$I=[\Lambda_\cap:N\Lambda_0]^r$, one has
\begin{equation}\label{eq:primary-index}
 v_\ell\!\left(|N_{E/\mathbb Q}(\det q)|^r\right)=v_\ell(I).
\end{equation}
If $I_\tau(L)$ denotes the fractional ideal generated by the determinants
of columns of $L$, transported by the coefficient embedding $\tau$,
then at every coefficient-field place $v$ above $\ell$,
\begin{equation}\label{eq:primary-det-ideal}
 v\!\left(\prod_\tau\tau(\det p)^{m_\tau}\right)
 =\sum_\tau m_\tau\bigl(v(I_\tau(\Lambda_\cap))
                         -v(I_\tau(\Lambda))\bigr).
\end{equation}
Thus the original ideal inclusion
$(s)\prod_\tau I_\tau(\Lambda_\cap)^{m_\tau}
 \subset\prod_\tau I_\tau(\Lambda)^{m_\tau}$ implies
$v(s\prod_\tau\tau(\det p)^{m_\tau})\geq0$.
\end{proposition}

\begin{proof}
At an unramified nonsplit finite place, a unimodular Hermitian lattice
has an orthonormal basis.  Indeed its reduction is a nondegenerate
Hermitian space over a quadratic finite field.  It has a vector of
nonzero norm, also in residue characteristic two: otherwise the
surjectivity of the residue trace, applied to
$h(x+ay,x+ay)$, would make the pairing zero.  Lift that vector.  Its
norm is a unit; the norm map on units of an unramified quadratic
extension is surjective, by the finite-field norm map and the
surjective trace on successive principal-unit quotients.  Rescale to
norm one, split off its orthogonal complement, and induct.  At a split
place the two integral factors are in perfect duality, and any basis
on one factor with its dual on the other gives the same conclusion.
This proves the assertion at the dyadic places as well.

Choose these local unitary identifications at the places in $S$.
They can be approximated simultaneously by a rational unitary matrix:
the Cayley map $X\mapsto(1-X)(1+X)^{-1}$ identifies the open subset
$\det(1+X)\ne0$ of the $F$-vector space $X^*=-X$ with the open subset
$\det(1+u)\ne0$ of the unitary group.  The latter meets every local open
neighbourhood.  Weak approximation in that vector space therefore
gives a $\gamma$ taking $\Lambda_{\mathfrak p}$ onto
$(\Lambda_0)_{\mathfrak p}$ for every $\mathfrak p\in S$; equality of
these lattices is an open condition.  Clearing the two inclusions away
from $S$ gives a positive $k_2$ prime to $2\ell$ with
$k_2\gamma\Lambda\subset\Lambda_0$ and
$k_2\Lambda_0\subset\gamma\Lambda$.

The projective module $\Lambda_\cap$ has rank $n$.  Chinese remainders
choose $n$ global elements that reduce to a basis at every prime in
$S$.  Let $g$ have these columns.  Nakayama's lemma makes them a basis
locally at $S$, and $\Lambda_\cap/g\Lambda_0$ is finite with no support
there.  Some positive integer $k_1$ prime to $2\ell$ annihilates it.
Define
\[
 k=k_1k_2,\qquad p=k_2\gamma g,\qquad q=g^{-1}k_1N I_n.
\]
The stated inclusions make $p,q$ integral and prove all the displayed
identities.  For $n=0$ use the empty matrices and $k=k_1=1$.

At a prime over $\ell$, $g\Lambda_0=\Lambda_\cap$ and $k_1$ is a unit.
Consequently $g$ identifies $\Lambda_0/q\Lambda_0$ with
$\Lambda_\cap/N\Lambda_0$ after localization.  The determinant formula
for the index of an integral matrix, or its Smith normal form over
each local discrete valuation ring, proves \eqref{eq:primary-index}.
Similarly $\gamma\Lambda=\Lambda_0$ locally and
$p=k_2\gamma g$, with $k_2$ a unit.  Taking determinant ideals yields
\eqref{eq:primary-det-ideal}, including ramified coefficient embeddings.
An inclusion of fractional ideals reverses the inequality on their
orders, giving the last assertion.  No principalization of either
global determinant ideal has been used.
\end{proof}

\subsection*{Finite height and splitting from the supplied model}

Let $R$ be a complete discrete valuation ring of characteristic zero,
with residue characteristic $\ell$ and fraction field $K$.  Suppose
the original smooth integral group model has semiabelian special fibre,
and its generic fibre is the prescribed abelian variety $A/K$.
Retain its coefficient action and its actual formal chart at the identity.
Every multiplication map $[a]$, $a\geq1$, has finite kernel on the
geometric special fibre: this holds for tori and abelian varieties and
passes to their extensions.  The completed local ring of this kernel
at the identity is therefore Artinian.  Equivalently, the coordinates
of $[a]$ have finite colength in the residue power-series ring.  This
supplies exactly the finite-height hypothesis of
Theorem~\ref{thm:formal-kernel-jacobian-trace}.

Suppose $s\mu=[cM]$, where $(c,\ell)=1$, and all geometric $M$-torsion
of $A$ is $K$-rational.  Multiplication by $c$ is a formal automorphism,
so the formal kernel of $\mu$ is killed by $M$.  Its coordinate ring is
finite flat over $R$, by finite height and the preceding formal trace
argument.  Its geometric generic points are integral and algebraize
inside the algebraic kernel: take an affine neighbourhood of the
identity and its completion; the finite quotient represents the same
formal closed subscheme over the complete base.  Such a point is
$K$-rational by the $M$-torsion hypothesis.  Normality of $R$ then puts
its coordinates in $R$.  More explicitly, the coordinates of the
affine neighbourhood belong to $K$ and are integral over $R$, and the
formal coordinates are their convergent $R$-adic limits.  This proves
the exact splitting condition, not just splitting after a field
extension.  Moreover a small point killed by $cM$ is $\ell$-primary,
since the prime-to-$\ell$ part acts invertibly on the formal group.

\subsection*{The full primary sum}

For the original symmetric rigidified line on $A$, write $T_u$ for its
canonical odd-level theta translation.  Write $H(P;F)$ for the original
harmonic value on an $n$-tuple of sections $F=(F_b)$; it is linear in
each column.  In this paragraph use the source's general factorization
$\mu=pq$, with $\mu^*\mu=N^2I_n$ and $N$ odd, and put
\[
 W=\ker(\mu)\cap A[\ell^\infty]^n,\quad
 W_0=\ker(q)\cap W,\quad
 W_1=W\cap\widehat A(\mathfrak m)^n,
 \qquad \delta_p=\prod_\tau\tau(\det p)^{m_\tau}.
\]
Assume the original rational-torsion condition, and assume
$gq=k_1N I_n$ with $(k_1,\ell)=1$.  Suppose the product Weil pairing is
one on $W\times W$.  Let each $F_b$ be fixed by the canonical translations
of $\ell$-primary $N$-torsion, and set $G_u=(T_{u_b}F_b)_b$.
For every $u\in W$, retain the regularity of $G_u$ on the original
integral chart, its original integral symmetric translation units, and
the chart bound $\nu(H)\geq-\rho$ for integral harmonic input, where
$\nu=v/e$ and $e>0$.  Finally assume $\lambda\mu_*P$ is integral in
that same harmonic chart, with the original differential normalization.

\begin{theorem}\label{thm:full-primary-theta-trace}
There is $z\in K$ with $\nu(z)\geq-\rho$ such that
\begin{equation}\label{eq:full-primary-theta-trace}
       \lambda\sum_{u\in W}H(P;G_u)=|W_0|\,\delta_p\,z.
\end{equation}
All of $W$ occurs in this identity, including the nonformal points.
There is no division of $z$ by the size of a torsion group.
\end{theorem}

\begin{proof}
For two odd-level canonical translates, their composition differs
from the canonical translate of the sum by a scalar $r_b$ with
$r_b^{N^2}=1$ and $r_b^2$ equal to the corresponding Weil pairing
(with the fixed source convention).  These identities follow by
applying respectively the order and inversion normalizations in the
theta group; equivalently, the odd-order canonical lift uses the
inverse of $2$ in $\mathbb Z/N^2\mathbb Z$.  Isotropy gives
$(\prod_b r_b)^2=1$, while $(\prod_b r_b)^{N^2}=1$.  Hence
$\prod_b r_b=1$.  Each $r_b$ is an integral unit, since it and its
inverse are roots of unity.  These are identities of the sections themselves, before applying any
harmonic functional.  On the external tensor product line over $A^n$,
the product action $T_u=\boxtimes_bT_{u_b}$ consequently satisfies
$T_uT_x=T_{u+x}$ for $u,x\in W$.  Multilinearity then makes the
same product scalar disappear from $H$.

If $u\in W_0$, then $k_1Nu=0$, and primary torsion cancels $k_1$.
Thus $Nu=0$.  The stipulated stabilizer property of $F$ and the preceding
composition calculation imply, for every $x\in W$, the stronger
section identity
\[
 T_u\Bigl(\boxtimes_bG_{x,b}\Bigr)
 =T_uT_x\Bigl(\boxtimes_bF_b\Bigr)
 =T_xT_u\Bigl(\boxtimes_bF_b\Bigr)
 =\boxtimes_bG_{x,b}.
\]
Equivalently the column factors relating the translated tuple to
$G_x$ are integral units with product one.  Restriction to the actual
formal chart, including its original theta gauge, preserves this
identity.  Thus the full product formal section is $W_0$-invariant;
this is the invariance required by formal trace.  It also implies
$H(P;G_{x+u})=H(P;G_x)$.
Set $f(x)=H(P;G_x)$ and $c=|W_0\cap W_1|$.  Apply the already proved
formal trace to the regular tuple $G_x$ and the full small kernel $W_1$.
The finite-height and splitting arguments above justify that this is
the entire formal kernel.  Its $q$-stabilizer is exactly $W_0\cap W_1$.
The section identity just established gives its required integral
product invariance, rather than only invariance of its harmonic value.  We obtain a $z_x$ satisfying
\[
 \nu(z_x)\geq-\rho,\qquad
       \lambda\sum_{u\in W_1}f(x+u)=c\,\delta_p z_x.
\]

Pass to $\overline W=W/W_0$ and let $J$ be the image of $W_1$ there.
Each element of $J$ has $c$ preimages.  The last identity is therefore
$c\lambda\sum_{j\in J}\bar f(\bar x+j)=c\delta_p z_x$.
Cancel the nonzero integer $c$ in $K$.  The surviving $z_x$ is the
same element with the same valuation bound.  Choose one representative
for each coset of $J$ in $\overline W$ and sum.  Finally each element of
$\overline W$ has $|W_0|$ preimages in $W$.  This gives
\eqref{eq:full-primary-theta-trace} with $z=\sum z_x$.
The ultrametric inequality gives $\nu(z)\geq-\rho$, also when the sum
vanishes.  In particular the argument does not require
$W_0\subset W_1$ or $W_1\subset W_0$.
\end{proof}

\subsection*{Exact cancellation in the original normalized sum}

Return to Proposition~\ref{prop:primary-adapted-construction} and the
original lattice-theta data.  Let the rational coefficient module of
$A$ have rank $r=\mathrm{p}+\mathrm{q}$ over $E$.  Each local Tate
module is free of rank $r$ over the corresponding factor of
$\mathcal O_E\otimes\mathbb Z_\ell$.  Smith normal form thus gives
\[
 |\ker(q)\cap A[\ell^\infty]^n|
 =\ell^{v_\ell(|N_{E/\mathbb Q}(\det q)|^r)}
 =\ell^{v_\ell(I)}=:I_\ell.
\]
The first assertion uses the original coefficient action, not the
complex tangent rank.  Since $gq=k_1NI_n$, these kernel points are
killed by $N$; the original rational $N^2$-torsion condition makes all
of them rational.  Also $\ker q\subset\ker(pq)$, so this is exactly
$|W_0|$.

The original lattice-annihilator definition of $W$ agrees with the
primary kernel of $pq$.  Indeed $k\gamma$ takes $\Lambda$ into
$\Lambda_0$, giving one direction.  Conversely, for any rational row
matrix $a$ taking $\Lambda$ into $\Lambda_0$, the inclusion
$k\Lambda_0\subset\gamma\Lambda$ makes
$k a\gamma^{-1}$ integral.  Multiplying its relation by $pq=kN\gamma$
gives $k^2Na$.  Cancellation of the prime-to-$\ell$ scalar $k^2$ on
primary torsion recovers every original row relation.  The adjoint
identity similarly recovers the $N^2$-torsion condition.  Thus no
different torsion subgroup has been substituted.

For $P=0$ the valuation assertion below is immediate.  Otherwise let
$d=\deg P$, $S_\ell$ be the original sum over this primary group,
and $\eta_k$ the exact chain-rule factor of the normalized pullback by
$[k]$.  Its original definition and the row-degree condition give
\[
 \eta_k=\operatorname{nrm}(k)
 =\prod_\tau\prod_{i=1}^{m_\tau}\tau(k)^{w(\tau,i)}
 =k^{\sum_{\tau,i}w(\tau,i)}=k^d.
\]
Indeed every monomial in the nonzero $P$ has row sums $w(\tau,i)$,
so their total sum is exactly $d$.  The original normalized pullback
chain rule supplies precisely $\operatorname{nrm}(k)$, with no extra
line scalar: the harmonic gauge transforms by the same rigidified
pullback, and each differential in row $(\tau,i)$ contributes
$\tau(k)$.  This is also the explicit source interface
\texttt{hv\_pulls}, together with
\texttt{nrm\_natCast'}/\texttt{totalDegree\_of\_rowDeg}.
We retain $\eta_k$ below to display every original normalization.  Pull back the original
sections by $[k]$.  Multiplication by $k$ permutes the primary group,
and $pq=kN\gamma$ gives the integral transformed polynomial
$(kN)^{-d}(pq)_*P=\gamma_*P$.  Its integrality follows from the
original coefficient lattice and $k\gamma\Lambda\subset\Lambda_0$,
because $k$ is a local unit.  Unitary $\gamma$ preserves harmonicity.
The required isotropy also concerns these exact objects.  On $A^n$ the
Rosati identity $(pq)^*(pq)=(kN)^2I_n$ identifies the polarisation
of $(pq)^*L_{\mathrm{box}}$ with that of
$L_{\mathrm{box}}^{(kN)^2}$.  The kernel of an isogeny is isotropic
for the commutator pairing of its pulled-back line: its descent lifts
commute.  Thus the product Weil pairing is one on the whole kernel.
The primary subgroup inherits this identity.  Functoriality of the
original normalized pullback transports the $N$-torsion stabilizer
to the $(kN)$-torsion stabilizer, since $k$ is invertible on primary
torsion; no new translation normalization is chosen.
The full primary trace now says, with no discarded scalar,
\[
       (kN)^{-d}\eta_k S_\ell=I_\ell\delta_p z,
                         \qquad \nu(z)\geq-\rho.
\]
Consequently
\begin{equation}\label{eq:primary-normalized-cancellation}
 s I^{-1}N^{-d}S_\ell
   ={I_\ell\over I}\,k^d\eta_k^{-1}\,(s\delta_p)z.
\end{equation}
The first three factors before $s\delta_p$ are local units.
The original determinant-ideal condition makes $s\delta_p$ integral
by \eqref{eq:primary-det-ideal}.  This proves the lower bound
$\nu(sI^{-1}N^{-d}S_\ell)\geq-\rho$ for the actual normalized primary
sum, conditional only on the stated geometric section inputs and the
original theta identities.  Summing the prime-to-$\ell$ cosets preserves
the same bound when their prescribed base tuples have those inputs.

The results apply to both manuscripts with their unchanged auxiliary
fields and ranks.  They supply the adapted arithmetic construction,
the full-primary passage and the determinant/index cancellation.
They do not construct the integral model or prove that all its reached
sections are regular.  That regularity, the prescribed mixed
level-torsion pullback identities, the raw theta norm and theta
spanning remain the next geometric obligations.  In particular the
global one-copy clearing integer is not declared proved here.
The exact source interfaces and their bounded hashes are recorded
separately\cite{primary-theta-source}.
